\sequence{Bilan de trimestre et remédiation}{Bilan de trimestre}{16--20 novembre 2026}
\section{Bilan individuel}
\begin{methode}[Consigne]
Chercher sans le cours pendant 25 minutes, puis vérifier avec les méthodes du livret. Pour chaque question, noter : réussi seul, réussi avec aide, ou à reprendre. Justifier les réponses.
\end{methode}
\ex[Bilan des séquences 1 à 4]{\begin{enumerate}
\item Rendre $90/126$ irréductible. Démontrer que la somme de deux entiers impairs est paire.
\item Avec $A(-1;2)$ et $B(3;-2)$, calculer $\vect{AB}$, $AB$ et le milieu de $[AB]$ dans un repère orthonormé.
\item Déterminer une équation de $(AB)$ ; $C(5;-4)$ appartient-il à cette droite ?
\item Calculer $[-2;4]\cap]1;7]$ et $[-2;4]\cup]1;7]$.
\item Résoudre $|x-1|\leq4$ et expliquer pourquoi $1/3\notin\D$.
\end{enumerate}}
\reponse[9]{1. $\pgcd(90,126)=18$, donc $90/126=5/7$. Si $u=2k+1$ et $v=2\ell+1$, alors $u+v=2(k+\ell+1)$ est pair.\\2. $\vect{AB}=(4;-4)$ ; $AB=4\sqrt2$ ; milieu $I(1;0)$.\\3. $y=-x+1$. Pour $C$, $-4=-5+1$ : il appartient à $(AB)$.\\4. Intersection $]1;4]$ ; réunion $[-2;7]$.\\5. $x\in[-3;5]$. Une écriture $1/3=a/10^n$ imposerait que $10^n$ soit divisible par $3$, ce qui est impossible.}
\prof{Barème possible sur 20 : 4 points par question. Repérer séparément maîtrise du calcul, choix de méthode et qualité de la justification.}
\fiche{Corriger une erreur et reprendre une méthode}
\ex[Analyse de productions]{Corriger : (a) « $\vect{AB}=(x_A-x_B;y_A-y_B)$ » ; (b) « Le vecteur nul dirige toutes les droites » ; (c) « $|x-2|\leq3$ donne $x\in[2;3]$ » ; (d) « $1/3\approx0{,}333$, donc $1/3$ est décimal ».}
\reponse[5]{(a) Soustraire origine à extrémité : $(x_B-x_A;y_B-y_A)$. (b) Un directeur doit être non nul. (c) La distance à $2$ est au plus $3$ : $x\in[-1;5]$. (d) $0{,}333$ est différent de $1/3$ : $3\times0{,}333=0{,}999\ne1$. Une approximation décimale ne détermine pas la nature du nombre approché.}
\ex[Parcours de reprise]{\begin{enumerate}
\item Réduire $72/108$, puis $\frac34-\frac16$.
\item Avec $P(1;1)$ et $Q(5;3)$, déterminer le milieu, un directeur et une équation de $(PQ)$.
\item Déterminer $]-\infty;3]\cap[-1;+\infty[$, puis résoudre $|x+2|<1$.
\end{enumerate}}
\reponse[5]{1. $2/3$ ; $7/12$. 2. Milieu $(3;2)$, directeur $(2;1)$, équation $x-2y+1=0$. 3. Intersection $[-1;3]$ ; solution de l'inéquation $]-3;-1[$.}
\fiche{Réinvestir et s'autoévaluer}
\ex[Problème de synthèse]{Deux déplacements sont représentés par $\vec u(6;4)$ et $\vec v(9;6)$.
\begin{enumerate}\item Montrer qu'ils sont colinéaires et trouver le coefficient reliant $\vec v$ à $\vec u$.
\item Le départ est $A(-2;1)$. Déterminer les arrivées $B$ et $C$ associées aux déplacements $\vec u$ et $\vec v$.
\item Trouver une équation de la droite passant par ces trois points.
\item Un point $M(x;y)$ de cette droite doit vérifier $|x-4|\leq3$. Donner l'intervalle des abscisses possibles.
\end{enumerate}}
\reponse[7]{1. $\vec v=\frac32\vec u$. 2. $B(4;5)$ et $C(7;7)$. 3. Un directeur est $(3;2)$ ; l'équation est $2x-3y+7=0$, vérifiée par les trois points. 4. $|x-4|\leq3\iff x\in[1;7]$.}
\begin{cours}[Mon bilan]
Ma priorité : \dotfill\par L'erreur à éviter : \dotfill\par La méthode à réutiliser : \dotfill\par Un exercice à refaire sans aide : \dotfill
\end{cours}
\prof{Valider la reprise par un exercice analogue réalisé sans aide après un délai. À l'oral, demander une phrase reliant calcul et conclusion géométrique.}
